% countdown-stp.tex
% Software Test Plan for the countdown program.
% Structure per IEEE Std 829-2008 / ISO/IEC/IEEE 29119-3:2021.
% Compiled with: pdflatex countdown-stp.tex

\documentclass[12pt]{article}

\usepackage{geometry}
\geometry{letterpaper, margin=1in}
\usepackage{hyperref}
\usepackage{booktabs}
\usepackage{array}
\usepackage{longtable}
\usepackage{enumitem}
\usepackage{titlesec}
\usepackage{fancyhdr}
\usepackage{xcolor}

\definecolor{cdteal}{RGB}{0,90,120}

\pagestyle{fancy}
\fancyhf{}
\renewcommand{\headrulewidth}{0.4pt}
\fancyhead[L]{\small Countdown Project}
\fancyhead[R]{\small STP: countdown}
\fancyfoot[C]{\thepage}

\titleformat{\section}{\large\bfseries\color{cdteal}}{\thesection}{0.5em}{}[\titlerule]
\titleformat{\subsection}{\normalsize\bfseries}{\thesubsection}{0.5em}{}

\newcommand{\code}[1]{\texttt{#1}}
\newcommand{\req}[1]{\texttt{#1}}

\begin{document}

\begin{titlepage}
  \centering
  \vspace*{2cm}
  {\Huge\bfseries\color{cdteal} Software Test Plan\\[0.5em]}
  {\Large countdown: An X~Window Countdown Timer}\\[2em]
  {\large Countdown Software Project}\\[0.5em]
  {\normalsize Prepared by: J.\,R.\,Evans}\\[2em]
  \begin{tabular}{ll}
    Document identifier: & CD-STP-001                                      \\
    Revision:            & 1.0                                             \\
    Date:                & June 13, 2026                                   \\
    Status:              & Released                                        \\
    Governing standard:  & ISO/IEC/IEEE 29119-3:2021 (Test Documentation)  \\
  \end{tabular}
  \vfill
  {\small\textit{This document is prepared in accordance with the
  principles of rational documentation articulated in
  D.\,L.\,Parnas and P.\,C.\,Clements, ``A Rational Design Process:
  How and Why to Fake It,'' \textit{IEEE Transactions on Software
  Engineering}, vol.\,SE-12, no.\,2, pp.\,251--257, February 1986.}}
\end{titlepage}

\tableofcontents
\clearpage

% =================================================================
\section{Introduction}
% =================================================================

\subsection{Document identifier}

CD-STP-001 Rev.\,1.0, \textit{Software Test Plan: countdown}.

\subsection{Purpose}

This Software Test Plan (STP) specifies how the requirements of CD-SRS-001
shall be verified by execution and static inspection.  It identifies the
test items, the test environment, the test cases and their pass/fail
criteria, the requirement-to-test traceability, and the deliverables.

\subsection{Scope}

The STP covers the \code{countdown} executable tangled from
\code{countdown.w}.  Tests divide into two classes: \emph{automated} cases,
run headless by the harness \code{test-countdown.sh} (driven by
\code{run-stp.sh}); and \emph{manual} cases, performed by an operator at a
live X~display where automation would require a windowing server and human
visual confirmation.

% =================================================================
\section{Test items}
% =================================================================

\begin{itemize}[noitemsep]
  \item \code{countdown.w} --- the literate source of record.
  \item \code{countdown.c} --- the tangled C source (build artefact).
  \item \code{countdown} --- the compiled executable.
  \item \code{Makefile} --- the build driver.
\end{itemize}

% =================================================================
\section{Features to be tested}
% =================================================================

All requirements in CD-SRS-001 \S3 are in scope: functional
(REQ-F-$\ast$), GUI (REQ-G-$\ast$), timing (REQ-T-$\ast$), design
constraint (REQ-D-$\ast$), and build (REQ-B-$\ast$).

% =================================================================
\section{Features not to be tested}
% =================================================================

The correctness of the X~server, the host's time-zone database, the C
compiler, and the CWEB tools themselves are outside scope; they are
assumed correct.

% =================================================================
\section{Test environment}
% =================================================================

\textbf{Automated.}  A POSIX host with \code{cc}, \code{make}, a Python~3
interpreter, and the CWEB tools (\code{ctangle}/\code{cweave}, reached
through \code{cweb.sh}).  No X~server is required: every automated case
either inspects source or exercises the program's pre-display code paths.

\medskip\noindent
\textbf{Manual.}  Additionally an X~display (Linux/X11 or macOS/XQuartz),
at which an operator can observe the window and operate the keyboard and
mouse.

% =================================================================
\section{Approach}
% =================================================================

Automated cases emit one line per requirement of the form
\code{PASS}/\code{FAIL}/\code{SKIP} followed by the requirement identifier
and a short description.  A case is \code{SKIP}ped only when its
precondition (e.g.\ a live X~display) is absent.  The run exits with a
non-zero status if any case fails.  Manual cases are recorded on the
checklist in Section~8.

% =================================================================
\section{Automated test cases}
% =================================================================

\begin{longtable}{p{0.12\textwidth}p{0.16\textwidth}p{0.62\textwidth}}
\toprule
\textbf{Case} & \textbf{Req.} & \textbf{Procedure / pass criterion} \\
\midrule
\endhead
TC-01 & REQ-B-01, REQ-D-04 & \code{make tangle} regenerates \code{countdown.c} from \code{countdown.w}; the executable builds. \\
TC-02 & REQ-D-05 & The compile produces no \code{warning:} lines under \code{-Wall}. \\
TC-03 & REQ-B-02 & The link line references \code{-lX11} and no widget toolkit. \\
TC-04 & REQ-D-01 & \code{countdown.w} contains CWEB markers (\code{@*}, \code{@<...@>=}) and names literate programming. \\
TC-05 & REQ-D-02 & No CWEB code section exceeds 23 lines of code (checked by script). \\
TC-06 & REQ-D-03 & A ``POSIX system services'' section exists and the index records each system call (\code{time}, \code{mktime}, \code{localtime}, \code{select}). \\
TC-07 & REQ-F-02 & The \code{DEFAULT\_TARGET} macro equals \code{2026-08-19 08:00:00}. \\
TC-08 & REQ-F-03 & The source accepts \code{argv[1]} as the target specification. \\
TC-09 & REQ-F-04 & Running \code{countdown "not-a-date"} prints a diagnostic and exits with status~2, opening no window. \\
TC-10 & REQ-F-05 & \code{fmt\_duration} emits the \code{d HH:MM:SS} form with zero-padding. \\
TC-11 & REQ-F-06 & \code{remaining\_secs} clamps a past target to zero (verified via a past-dated target with status~2 guard / source inspection of the clamp). \\
TC-12 & REQ-T-02, REQ-T-03 & The source computes remaining time with \code{time} and converts the target with \code{mktime} using \code{tm\_isdst = -1}. \\
TC-13 & REQ-T-04 & The event loop blocks on \code{select} (no busy-wait). \\
\bottomrule
\end{longtable}

% =================================================================
\section{Manual test cases}
% =================================================================

\begin{longtable}{p{0.12\textwidth}p{0.16\textwidth}p{0.62\textwidth}}
\toprule
\textbf{Case} & \textbf{Req.} & \textbf{Procedure / pass criterion} \\
\midrule
\endhead
MC-01 & REQ-G-01, REQ-G-02 & Launch \code{countdown}; a window opens showing target, remaining time, status, and hint. \\
MC-02 & REQ-F-01, REQ-T-01 & With no argument, the window counts down toward the default target, the seconds advancing once per second. \\
MC-03 & REQ-G-03 & Press space (or click): the status toggles running$\leftrightarrow$paused and the figure freezes when paused. \\
MC-04 & REQ-G-04 & Press \code{q} or Escape: the window closes and the program exits with status~0. \\
MC-05 & REQ-G-05 & Resize and re-expose the window: contents are repainted correctly. \\
MC-06 & REQ-F-07 & Launch with a target a few seconds ahead; on arrival the status shows the arrival message and the figure reads \code{0 d 00:00:00}. \\
\bottomrule
\end{longtable}

% =================================================================
\section{Pass/fail criteria}
% =================================================================

The test campaign passes when every automated case reports \code{PASS}
(or a justified \code{SKIP} for an absent precondition) and every manual
case is confirmed by the operator.  Any \code{FAIL} fails the campaign and
blocks release.

% =================================================================
\section{Traceability matrix}
% =================================================================

\begin{longtable}{p{0.20\textwidth}p{0.30\textwidth}p{0.40\textwidth}}
\toprule
\textbf{Requirement} & \textbf{Verified by} & \textbf{Method} \\
\midrule
REQ-F-01 & MC-02 & Manual \\
REQ-F-02 & TC-07 & Automated (static) \\
REQ-F-03 & TC-08 & Automated (static) \\
REQ-F-04 & TC-09 & Automated (dynamic) \\
REQ-F-05 & TC-10 & Automated (static) \\
REQ-F-06 & TC-11 & Automated (static) \\
REQ-F-07 & MC-06 & Manual \\
REQ-G-01 & TC-03, MC-01 & Automated + Manual \\
REQ-G-02 & MC-01 & Manual \\
REQ-G-03 & MC-03 & Manual \\
REQ-G-04 & MC-04 & Manual \\
REQ-G-05 & MC-05 & Manual \\
REQ-G-06 & TC-04 (fallback code present) & Automated (static) \\
REQ-T-01 & MC-02 & Manual \\
REQ-T-02 & TC-12 & Automated (static) \\
REQ-T-03 & TC-12 & Automated (static) \\
REQ-T-04 & TC-13 & Automated (static) \\
REQ-D-01 & TC-04 & Automated (static) \\
REQ-D-02 & TC-05 & Automated (static) \\
REQ-D-03 & TC-06 & Automated (static) \\
REQ-D-04 & TC-01 & Automated (dynamic) \\
REQ-D-05 & TC-02 & Automated (dynamic) \\
REQ-B-01 & TC-01 & Automated (dynamic) \\
REQ-B-02 & TC-03 & Automated (static) \\
\bottomrule
\end{longtable}

% =================================================================
\section{Test deliverables}
% =================================================================

\begin{itemize}[noitemsep]
  \item This plan (CD-STP-001).
  \item The automated harness \code{test-countdown.sh}.
  \item The plan runner \code{run-stp.sh}.
  \item The harness console log and exit status from each run.
\end{itemize}

% =================================================================
\section{Responsibilities and schedule}
% =================================================================

The maintainer runs the automated harness before every release; a \code{0}
exit status is a precondition for tagging.  The manual checklist
(Section~8) is performed at a live display whenever GUI-affecting code
changes.

\end{document}
